\chapter{Ekstraksjonen}

\begin{motto}
Tala er ikkje biverknader. Dei er primæreffekten av eit designval.\\ Dei som tok valet, heiter designarar.
\end{motto}

Vi har følgt ei line av skade gjennom heile historia: den planlagde foreldinga som amerikansk industridesign gjorde til forretningsmodell, eingongskulturen, mengda gjenstandar vi eig og kastar. I dette kapitlet når lina si tydelegaste form. Dei selskapa som i dag er mellom dei høgast verdsette i verda, byggjer forretningsmodellen sin på å ekstrahere privat informasjon, fange merksemd og produsere avhengigheit. Dette er ikkje misbruk av design. Det er design, utført kompetent, langs den eine aksen oppdraget spesifiserte, av folk med designgrader, premiert på designkonferansar. At faget ikkje har noko internt språk for å kalle det noko anna enn vellukka, er den skarpaste målinga av kva eit fundament ville gjort, og ikkje gjer.

\section{Frå brukar til råvare}

Skiftet har eit namn. Shoshana Zuboff kalla det \emph{overvakingskapitalisme}: ein økonomisk logikk som einsidig gjer menneskeleg erfaring om til åtferdsdata, og desse data om til prediksjonsprodukt som vert selde i marknader for framtidig åtferd \autocite{zuboff2019}. I denne logikken er ikkje brukaren kunden. Brukaren er råvara. Førre kapittel om brukarsentrering får her si mørke avslutning: den same metodikken som lærte faget å studere brukaren empatisk, vart instrumentet for å kjenne henne godt nok til å fange henne. Tim Wu spora den lengre historia — frå pressa via radio og fjernsyn til strøymen — og viste at «merksemdshandlarane» alltid har levd av å selje publikum til annonsørar, men at digitale grensesnitt gjorde fangsten kontinuerleg og personleg \autocite{wu2016}.

Det avgjerande er at fangsten er \emph{designa}. Den skjer ikkje av seg sjølv; ho er bygd inn i forma — i varslane, i den uendelege rullinga, i dei raude prikkane, i belønningsplanane som etterliknar spelautomatens. Nick Seaver har vist korleis tilrådingssystem best forståast som \emph{feller}: strukturar designa for å fange og halde, lånte direkte frå fangstteknologiens logikk \autocite{seaver2019captivating}. Forma er ikkje nøytral og så misbrukt; forma \emph{er} fangstmekanismen.

\section{Avhengigheit som designmål}

At dette er eit eksplisitt designmål, og ikkje ein utilsikta følgje, kan dokumenterast frå faglitteraturen til feltet sjølv. B.J. Fogg etablerte ved Stanford eit «Persuasive Technology Lab» og ein heil disiplin — captology — for korleis datamaskiner kan endre kva menneske tenkjer og gjer \autocite{fogg2003}. Nir Eyal gjorde innsikta om til ei praktisk oppskrift i \textit{Hooked}: korleis byggje vaneskapande produkt gjennom ein sløyfe av utløysar, handling, variabel belønning og investering \autocite{eyal2014hooked}. Adam Alter dokumenterte den andre sida — korleis desse same teknikkane produserer reell, klinisk åtferdsavhengigheit \autocite{alter2017}. Dette er ikkje kritikarar som skuldar bransjen for noko han nektar for; Fogg og Eyal \emph{lærer bort} mekanismen som ein kompetanse. Designaren av eit slikt system står framfor nett det same valet som designaren av eingongskoppen: kor mange aksar skal forma bere. Og fordi faget manglar eit fundament som tvingar fram fleire enn den eine kunden betalar for — engasjement, retensjon, konvertering — er det fullt mogleg å gjere ein strålande fagleg jobb medan ein byggjer ei felle.

\section{Signaturen i tala}

Skaden er ikkje abstrakt, og han let seg måle. James Williams, sjølv tidlegare Google-tilsett, argumenterte i \textit{Stand Out of Our Light} for at merksemdsøkonomien systematisk underminerer nett dei evnene vi treng for å styre liva våre \autocite{williams2018}. Jonathan Haidt har samla dokumentasjonen for kva som skjedde med ein heil generasjon då smarttelefonen og dei sosiale plattformene traff barndomen: kraftige auke i angst, depresjon og sjølvskade hjå unge, tett korrelert i tid med omlegginga av barndomen rundt skjermen \autocite{haidt2024}. Jean Twenge dokumenterte det same generasjonsskiftet i åtferd og psykisk helse \autocite{twenge2017}. Og den breiare litteraturen om algoritmisk skade — Cathy O'Neils «weapons of math destruction» \autocite{oneil2016}, Safiya Nobles påvising av korleis søkjemotorar forsterkar rasisme \autocite{noble2018} — viser at mønsteret ikkje er avgrensa til merksemd: overalt der eit system optimaliserer éin målbar akse og er blind for resten, ber forma signaturen av blindskapen.

Poenget for denne boka er ikkje å fordømme den einskilde designaren. Eit individuelt nei tapar berre oppdraget; det neste byrået tek det. Poenget er at skaden er \emph{kollektiv og strukturell}, og difor berre kan kurerast kollektivt og strukturelt — av eit fag som hadde ein standard å utelukke for brot på. At skulda er kollektiv, er nett kvifor ho er så lett å bere for kvar einskild: alle gjekk på dei same konferansane, vann dei same prisane, følgde den same prosessen. Ingen braut ein regel, fordi det ikkje fanst nokon.

\vignett

Når objektet er merksemda til eit barn, er forma framleis berre optimert for éin akse. I neste kapittel er objektet eit menneskeliv, og aksen er statleg vald. Det er den ytste konsekvensen av eit fag som ikkje kan seie nei.

\vidarelesing{\fullcite{zuboff2019}; \fullcite{wu2016}; \fullcite{fogg2003}; \fullcite{eyal2014hooked}; \fullcite{haidt2024}; \fullcite{seaver2019captivating}.}
